\chapter{진공 에너지와 드시터 지평선}

\lectureoverview{진공 에너지는 보통 물질과 다르다. 진공 에너지는 특별한 관성계를 선택하지 않으며, 우주가 팽창해도 희석되지 않는다. 중력을 포함하면, 겉보기에는 무해한 이 에너지의 영점이 시공간의 기하를 바꾼다. 물질 지배 법칙 \(a(t)\propto t^{2/3}\)을 유도하고, 이를 진공 지배 법칙 \(a(t)\propto e^{Ht}\)과 비교한 다음, 드시터 시공간을 등각 좌표로 다시 표현하겠다. 이렇게 얻은 인과 도표에는 고유 거리가 일정한, 관측자에 따라 달라지는 사건 지평선이 드러난다.}

\section{진공의 양자 요동}

강의 전체의 출발점이 되는 질문부터 살펴보자. 양자장론은 빈 공간에도 에너지가 있다고 말한다. 그렇다면 실제 전자--양전자 쌍이 무에서 끊임없이 생겨난다는 뜻일까?

\begin{classroomqa}
\audiencequestion{전자--양전자 쌍의 생성과 소멸은 진공에서 실제로 일어나는가, 아니면 단지 계산을 위한 장치인가?\lecturetimestamp{00:00:06}}
\lecturerresponse{가상 쌍은 진공 에너지에 기여하지만, 빈 공간에서 나타나는 모습을 관찰할 수 있는 보통 입자는 아니다. 이들은 양자 바닥 상태의 요동이며, 가상 과정에 적용되는 한계 안에서만 지속된다.}
\end{classroomqa}

조화 진동자는 가장 간명한 모형이다. 고전적으로 진동자는 퍼텐셜의 최솟값에서 정지해 있을 수 있다. 양자역학적으로는 그 최솟값에서 위치가 정확히 정해지는 것과 운동량이 정확히 0인 것이 양립할 수 없다. 바닥 상태는 두 조건 사이에서 타협하며 영점 에너지를 가진다.
\begin{equation}
E_0=\frac{1}{2}\hbar\omega.
\label{eq:l08-zero-point}
\end{equation}
이것은 진동자가 가질 수 있는 최소 에너지다. 들뜬 상태가 아니며, 관측 가능한 입자들이 흘러나오는 것을 뜻하지도 않는다.

자유장은 정상 모드들로 분해되며, 각 모드는 진동자처럼 거동한다. 공동 안의 전자기장이 익숙한 예다. 각 모드는 식~\eqref{eq:l08-zero-point}에 나오는 형태의 항을 보탠다. 전자장, 쿼크장 및 다른 양자장도 저마다 바닥 상태 요동에 기여한다. 이 기여들을 합한 것이 진공 에너지다. 순진한 미시적 추정값과 관측된 우주론적 척도 사이의 악명 높은 불일치는 실제 문제지만, 이 강의의 주제는 아니다.

\begin{classroomqa}
\audiencequestion{대략 \(120\)자릿수에 이르는 이 차이를 다시 논의하나요?\lecturetimestamp{00:02:59}}
\lecturerresponse{여기서는 다루지 않습니다. 지금의 목표는 우주 지평선을 이해하는 것입니다. 우주론 방정식이 어떻게 가속 팽창을 낳는지, 그리고 가속 팽창이 왜 사건 지평선을 만드는지를 살펴보겠습니다.}
\end{classroomqa}

\begin{classroomqa}
\audiencequestion{지난 강의에서 다룬 자외선--적외선 관계도 오늘 밤 논의하나요?\lecturetimestamp{00:03:59}}
\lecturerresponse{나중에 다시 다룰 수도 있지만, 예정된 논의에는 포함되지 않습니다. 우선 진공 에너지에서 우주론으로, 이어서 지평선으로 나아가겠습니다.}
\end{classroomqa}

\section{진공 에너지는 왜 로런츠 불변인가}

다음 질문은 단순해 보이지만 진공 에너지의 정의적 성질을 드러낸다. 빈 공간에 양의 에너지 밀도가 있다면, 왜 광학 매질처럼 빛의 속도를 바꾸지 않을까?

\begin{classroomqa}
\audiencequestion{양의 진공 에너지 밀도는 굴절 효과를 일으켜 빛을 조금 느리게 이동하게 해야 하지 않나요?\lecturetimestamp{00:04:35}}
\lecturerresponse{보통 물질에서는 그럴 수 있습니다. 물질이 특별한 정지계를 정하기 때문입니다. 진공 에너지는 다릅니다. 그런 계를 선택하지 않으므로, 모든 관성 관측자는 같은 진공 에너지 밀도를 측정합니다. 따라서 빛의 국소 측정값은 불변 속도 \(c\)를 유지합니다.}
\end{classroomqa}

법칙의 대칭성과 특정 상태의 대칭성을 구별해야 한다. 물리 법칙의 방정식은 병진 불변일 수 있지만, 스탠퍼드가 버클리와 같지 않고 지구가 성간 공간과 같지 않듯이 특정한 물질 배치는 그렇지 않을 수 있다. 물리적 배치 \emph{전체}를 병진 이동해도 법칙은 변하지 않는다. 그러나 실제 물질 배치가 모든 위치에서 똑같아 보일 필요는 없다. 마찬가지로, 방정식이 회전 불변이라고 해서 이 강의실이 회전 대칭인 것은 아니다.

로런츠 변환에도 같은 구별이 적용된다. 기체에는 정지계가 있다. 그 안을 움직이는 관측자는 분자들이 움직이는 것으로 보고, 다른 에너지 밀도를 측정한다. 따라서 기체의 지배 방정식은 로런츠 불변이어도 기체의 상태는 로런츠 불변이 아니다. 보통 물질, 복사 열욕, 실험 장치는 모두 특별한 운동이나 방향을 정해 대칭성을 깨뜨린다.

진공 에너지는 이런 일을 하지 않는 예외적인 형태의 물질이다. 진공 에너지 밀도를 측정하는 검출기를 생각해 보자. 검출기가 정지해 있을 때 어떤 값 \(\rho_{\rm vac}\)을 측정한다. 검출기를 가속한 다음 다시 측정해도 결과는 변하지 않는다. 상대론 표기에서는 이 성질을 계량 텐서에 비례하는 응력 텐서로 나타낸다.
\begin{equation}
T_{\mu\nu}^{\rm vac}=-\rho_{\rm vac}g_{\mu\nu},
\qquad
p_{\rm vac}=-\rho_{\rm vac}.
\label{eq:l08-vacuum-stress}
\end{equation}
\editorialnote{압력 관계는 판서에서 유도하지 않는다. 이 식은 로런츠 불변성과 팽창 중 일정한 밀도라는 강의의 두 정의적 성질을 뒷받침하는 공변적 진술이다.}

국소적이라는 단서가 중요하다. 관측자를 스쳐 지나가는 빛은 가까운 시계와 자로 측정할 때 \(c\)의 속도로 움직인다. 팽창하는 기하에서 멀리 떨어진 빛의 속도를 우리에 대해 묻는다면, 좌표계를 선택하고 멀리 떨어진 사건들을 비교하는 방법을 정해야 한다. 그 좌표 속도는 \(c\)와 같지 않을 수 있지만, 모든 국소 관측자는 \(c\)를 측정한다.

\begin{classroomqa}
\audiencequestion{이 논의에서 진공 에너지와 우주 상수는 같은 뜻인가요?\lecturetimestamp{00:11:57}}
\lecturerresponse{그렇습니다. 아인슈타인 방정식에서 같은 항을 놓는 방식이 다를 뿐입니다. 기하 쪽에서는 우주 상수로, 원천 쪽에서는 로런츠 불변 진공 응력 에너지로 표현할 수 있습니다.}
\end{classroomqa}

\section{희석되는 것과 그렇지 않은 것}

질량 \(m\)을 가진 비상대론적 입자들로 큰 영역을 채우자. 은하와 은하단을 평균 낼 만큼 큰 척도에서는 분포가 균일하다고 가정하자. 공변 부피 안에 입자 \(N\)개가 들어 있다면 물질 밀도는 다음과 같다.
\begin{equation}
\rho_m=\frac{Nm}{V}.
\end{equation}
모든 선형 치수가 두 배가 되면 부피는 \(2^3=8\)배 커지는 한편, \(N\)와 \(m\)는 변하지 않는다. 따라서 밀도는 이전 값의 8분의 1로 떨어진다. 같은 입자들이 더 큰 물리적 부피를 차지하므로 보통 물질은 희석된다.

복사도 희석되며, 그 속도는 더 빠르다. 광자 수 밀도는 부피가 커짐에 따라 낮아지고, 각 광자의 파장은 늘어나며 에너지는 적색편이된다. 뒤에서 결과 \(\rho_r\propto a^{-4}\)를 적겠지만, 오늘 밤 복사 지배 우주론을 풀지는 않겠다.

진공 에너지는 예외다.
\begin{equation}
\rho_{\rm vac}(t)=\text{constant}.
\label{eq:l08-rho-vac-constant}
\end{equation}
진공 에너지는 뭉치지 않고, 방향이나 관성계를 선택하지 않으며, 공간이 팽창해도 밀도가 낮아지지 않는다. 공변 상자가 커지면 그 안에 포함된 진공 에너지는 부피에 비례해 증가한다.

\editorialnote{일정한 밀도라고 해서 팽창하는 공변 상자 안의 진공 에너지 총량이 일정하다는 뜻은 아니다. 공변 응력 에너지 보존 법칙,
\(\dot\rho+3H(\rho+p)=0\)은 진공 에너지가 음의 압력 \(p_{\rm vac}=-\rho_{\rm vac}\)을 가지므로 성립한다. 일 항이 상자 안 진공 에너지 총량의 변화를 정확히 공급한다.}

\section{중력은 에너지의 영점을 본다}

중력을 고려하지 않는 실험실 물리에서는 모든 에너지에 같은 상수를 더해도 대개 아무것도 달라지지 않는다. 스펙트럼은 차이에 따라 정해진다. 들뜬 원자와 들뜨지 않은 원자 사이, 또는 입자가 있는 상태와 없는 상태 사이의 에너지가 그 예다. 빈 공간에는 절대 에너지를 잴 기준이 없어 보인다.

중력이 문제를 바꾼다. 질량은 에너지이고, 에너지는 중력의 원천이므로 진공 에너지도 중력을 만든다. 원자 전이에서는 보이지 않던 상수가 시공간을 휘게 할 수 있다.

아인슈타인 방정식은 기하와 원천을 구분한다.
\begin{equation}
G_{\mu\nu}=8\pi G_N T_{\mu\nu},
\label{eq:l08-einstein}
\end{equation}
\(c=1\)인 단위계에서 아인슈타인 텐서 \(G_{\mu\nu}\)는 계량 \(g_{\mu\nu}(x)\)과 그 도함수로 이루어진다. 응력 텐서 \(T_{\mu\nu}\)에는 물질과 비중력장의 에너지 밀도, 운동량 밀도, 압력 및 응력이 들어 있다. 진공 에너지도 보통 물질과 마찬가지로 이 원천 쪽에 놓인다.\editorialnote{강의의 구두 도입에서는 \(8\pi/3\)을 텐서 방정식과 잠시 연결해 말한다. 계수 \(1/3\)는 아래의 공간적으로 평탄한 프리드만 방정식에 들어가며, 식~\eqref{eq:l08-einstein}에는 들어가지 않는다.}

\begin{figure}[htbp]
\centering
\includegraphics[width=0.62\textwidth]{lecture_08_figure_02.png}
\caption{시공간 계량 \(g_{\mu\nu}(x)\)으로 이루어진 아인슈타인 방정식의 기하 쪽.}
\lectureframe{00:21:58.820}
\end{figure}

\section{척도 인자로 줄인 기하}

이제 공간 단면이 평탄한 가장 단순한 균일 팽창 기하를 택하자.
\begin{equation}
ds^2=-dt^2+a(t)^2\,d\mathbf{x}^2,
\qquad
d\mathbf{x}^2=dx^2+dy^2+dz^2.
\label{eq:l08-frw-metric}
\end{equation}
좌표 표식 \(\mathbf{x}\)는 무한한 고무판에 그려져 있다. 표식은 움직이지 않아도 되지만, 표식 사이의 물리적 거리는 달라진다. \(t\)가 일정할 때 \(\Delta\mathbf{x}\)만큼 떨어진 두 공변 점 사이의 고유 거리는 다음과 같다.
\begin{equation}
D(t)=a(t)\lvert\Delta\mathbf{x}\rvert.
\label{eq:l08-proper-distance}
\end{equation}
따라서 이처럼 대칭성이 높은 기하의 시간 의존성은 하나의 함수, 즉 척도 인자 \(a(t)\)에 모두 담긴다. \(a\)가 일정하면 공간을 재조정한 뒤 평탄한 민코프스키 시공간을 얻는다. \(a\)가 시간에 따라 변하면 각 공간 단면은 평탄한 채로 전체 시공간은 휘어질 수 있다.

\begin{classroomqa}
\audiencequestion{좌표 표식 \(x=7\)과 \(x=13\) 사이의 좌표 간격은 \(4\)인가요?\lecturetimestamp{00:27:40}}
\lecturerresponse{아닙니다. 간격은 \(13-7=6\)이므로 물리적 거리는 \(6a(t)\)입니다. 이 정정은 유용합니다. \(x\)는 공변 위치를 나타내고, 여기에 \(a(t)\)를 곱해야 물리적 거리가 됩니다.}
\end{classroomqa}

이 계량에서 에너지 밀도가 공간적으로 균일하면 아인슈타인 방정식의 한 성분은 공간적으로 평탄한 프리드만 방정식이 된다.
\begin{equation}
\left(\frac{\dot a}{a}\right)^2
=\frac{8\pi G_N}{3}\rho(t).
\label{eq:l08-friedmann}
\end{equation}
비율
\begin{equation}
H(t)\equiv\frac{\dot a(t)}{a(t)}
\label{eq:l08-hubble-definition}
\end{equation}을 허블 매개변수라고 한다. 이 단계에서는 정의일 뿐이며, 시간에 따라 달라질 수 있다.

\begin{figure}[htbp]
\centering
\includegraphics[width=0.90\textwidth]{lecture_08_figure_03.png}
\caption{판서에서 누락된 인자 \(1/a^2\)를 복원한 뒤의 평탄 팽창 계량과 프리드만 방정식.}
\lectureframe{00:30:30.000}
\end{figure}

\paragraph{에너지 균형 그림.}
이 에너지 균형 그림에서는 식~\eqref{eq:l08-friedmann}을 개략적으로 다음과 같이 쓴다.
\begin{equation}
-\frac{3}{8\pi G_N}H^2+\rho=0.
\label{eq:l08-energy-heuristic}
\end{equation}
물질 에너지는 양수인 반면, 팽창하는 척도 인자로 나타낸 중력 자유도는 음의 운동 에너지 항을 가진다. 이 축약된 문제에서는 음의 중력 에너지와 양의 물질 에너지의 합이 0이라는 것이 유용한 기억법이다.

\editorialnote{식~\eqref{eq:l08-energy-heuristic}은 대칭성이 매우 높은 미니초공간 모형의 해밀턴 제약식이다. 임의의 모든 우주에 좌표와 무관한, 전역적으로 보존되는 총에너지가 존재하며 그 값이 0이라는 정리는 아니다. 강의에서는 프리드만 제약식의 직관적 해석으로 이를 사용한다.}

\section{물질 지배 팽창}

한 변의 좌표 길이가 \(\Delta x=1\)인 공변 수학적 상자로 돌아가자. 물리적 한 변의 길이는 \(a(t)\)이므로 물리적 부피는 \(a(t)^3\)다. 상자에 질량 \(m\)인 비상대론적 입자 \(N\)개가 들어 있다면,
\begin{equation}
\rho_m(t)=\frac{Nm}{a(t)^3}
=\frac{\rho_0}{a(t)^3}.
\label{eq:l08-matter-density}
\end{equation}
입자 자체가 팽창하는 것은 아니다. 평균 간격이 팽창한다.

\begin{figure}[htbp]
\centering
\includegraphics[width=0.84\textwidth]{lecture_08_figure_04.png}
\caption{공변 상자 논증을 구성하며 제시한 계량과 프리드만 방정식.}
\lectureframe{00:36:07.940}
\end{figure}

식~\eqref{eq:l08-friedmann}에 대입하면 다음을 얻는다.
\begin{equation}
\left(\frac{\dot a}{a}\right)^2
=\frac{8\pi G_N\rho_0}{3a^3}.
\end{equation}
\(K=8\pi G_N\rho_0/3\)로 놓고 팽창하는 해를 택하자.
\begin{align}
\dot a^2&=\frac{K}{a},\\
\dot a&=\sqrt{K}\,a^{-1/2},\\
a^{1/2}\,da&=\sqrt{K}\,dt.
\label{eq:l08-matter-separated}
\end{align}

\begin{classroomqa}
\audiencequestion{변수를 분리한 방정식은 어떻게 풀어야 하나요?\lecturetimestamp{00:38:15}}
\lecturerresponse{양변을 적분하면 됩니다. \(a^{1/2}\,da\)의 적분은 \(\frac{2}{3}a^{3/2}\)이지 \(\frac{3}{2}a^{3/2}\)가 아닙니다. \(\frac{3}{2}\)은 새로운 지수이고, \(\frac{2}{3}\)는 계수입니다.}
\end{classroomqa}

따라서
\begin{equation}
\frac{2}{3}a^{3/2}
=\sqrt{K}\,(t-t_B),
\end{equation}
여기서 \(t_B\)는 선택한 시간 원점을 정하는 적분 상수다. 정규화와 시간 이동을 제외하면,
\begin{equation}
a(t)\propto (t-t_B)^{2/3}.
\label{eq:l08-matter-solution}
\end{equation}

\begin{figure}[htbp]
\centering
\includegraphics[width=0.86\textwidth]{lecture_08_figure_05.png}
\caption{물질 밀도를 대입하고 변수를 분리해 \(a^{1/2}da\propto dt\)에 이르는 과정.}
\lectureframe{00:38:30.000}
\end{figure}

곡선 \(t^{2/3}\)은 직선보다 느리게 증가한다. 기울기가 감소하므로 물질 지배 우주는 감속한다. 강의에서 사용한 뉴턴 역학의 표현으로 말하면, 입자들은 처음에 서로 멀어지고 있었고, 입자 사이의 중력이 그 분리 속도를 계속 낮춘다.

\begin{figure}[htbp]
\centering
\includegraphics[width=0.86\textwidth]{lecture_08_figure_06.png}
\caption{직선 팽창과 아래로 오목한 물질 지배 법칙 \(a(t)\propto t^{2/3}\)의 판서상 비교.}
\lectureframe{00:40:00.000}
\end{figure}

비교를 위해 복사 열욕의 경우는 다음과 같다.
\begin{equation}
\rho_r\propto a^{-4},
\end{equation}
\(a^{-3}\)와 달리, 부피 희석으로 세제곱 인자가 생기고 우주론적 적색편이로 인자가 하나 더 생긴다. 이에 해당하는 척도 인자도 풀 수 있지만, 강의에서는 그 계산을 일부러 다루지 않는다.

\section{진공 지배 팽창}

이제 프리드만 방정식의 우변에 진공 에너지만 놓자. \(\rho_{\rm vac}\)가 일정하므로,
\begin{equation}
\left(\frac{\dot a}{a}\right)^2
=\frac{8\pi G_N}{3}\rho_{\rm vac}
\equiv H^2.
\label{eq:l08-de-sitter-friedmann}
\end{equation}
구별해야 할 \(H\)의 용법이 두 가지 있다. 식~\eqref{eq:l08-hubble-definition}은 시간에 따라 달라질 수 있는 양 \(\dot a/a\)을 정의한다. 식~\eqref{eq:l08-de-sitter-friedmann}은 순수한 진공 에너지 우주에서 그 값이 다음 상수임을 말한다.
\begin{equation}
H=\sqrt{\frac{8\pi G_N\rho_{\rm vac}}{3}}.
\end{equation}
팽창하는 해를 택하면
\begin{equation}
\dot a=Ha,
\end{equation}
따라서
\begin{equation}
a(t)=C e^{Ht}.
\label{eq:l08-exponential-solution}
\end{equation}
적분 상수 \(C\)는 공변 좌표의 임의적 정규화를 정한다. 현재 \(a\)의 값에 비례하는 증가는 지수적 증가다.

\begin{figure}[htbp]
\centering
\includegraphics[width=0.88\textwidth]{lecture_08_figure_07.png}
\caption{진공 지배 방정식 \((\dot a/a)^2=H^2\), 그 1계 형태 \(\dot a=Ha\), 그리고 지수 해.}
\lectureframe{00:47:45.000}
\end{figure}

공변 물체에 대한 허블 법칙은 다음과 같다.
\begin{equation}
v_{\rm rec}=HD.
\label{eq:l08-hubble-law}
\end{equation}
이 순수한 드시터 사례에서 후퇴 속도가 \(c\)와 같아지는 고유 거리는 일정하다.
\begin{equation}
D_H=\frac{c}{H},
\label{eq:l08-hubble-radius}
\end{equation}
또는 판서에서 사용한 \(c=1\) 단위계에서는 \(D_H=H^{-1}\)이다. 이 고정된 척도는 기하학적으로 드시터 사건 지평선의 반지름으로 다시 나타난다.\editorialnote{허블 반지름이 임의의 시간 의존 우주론에서 자동으로 사건 지평선이 되는 것은 아니다. 정확한 드시터 공간에서 \(H\)가 일정하기 때문에 여기서는 두 개념이 일치한다.}

강의에서는 진공 에너지 검출기를 생각해 보는 사고 실험을 제시한다. 다른 원천에서 멀리 떨어진 곳에 매우 가벼운 시험 물체 두 개를 두고, 처음에는 공변 흐름에 대해 정지시킨 다음, 둘 사이의 중력은 무시할 수 있다고 하자. 두 물체의 간격은 지수적으로 팽창한다. 이 효과는 지상 장치로 측정하기에는 너무 작지만, 우주적 거리를 두고 과거를 되돌아보면 같은 물리를 관측할 수 있다.

\section{물질 지배에서 진공 지배로}

강의에서 다루는 실제 우주론 모형에는 두 항이 모두 들어 있다.
\begin{equation}
H(t)^2
=\frac{8\pi G_N}{3}
\left(\frac{\rho_0}{a(t)^3}+\rho_{\rm vac}\right).
\label{eq:l08-mixed-friedmann}
\end{equation}
\(a\)가 작을 때는 \(a^{-3}\)이 크므로 물질 항이 지배적이고, 해는 \(t^{2/3}\)와 비슷하다. 우주가 커지면 물질은 희석되는 반면 진공 에너지는 일정하게 유지된다. 결국 상수항이 우세해지고 척도 인자는 위쪽으로 휘며 \(e^{Ht}\)에 가까워진다.

\begin{figure}[htbp]
\centering
\includegraphics[width=0.87\textwidth]{lecture_08_figure_08.png}
\caption{강의에서 그린 정성적 전이: 초기의 물질 지배 감속에서 후기의 진공 지배 지수 증가로 이어진다.}
\lectureframe{00:53:30.000}
\end{figure}

\begin{classroomqa}
\audiencequestion{지수 해에는 적분 상수가 들어갑니다. 그 값이 극도로 작지 않다는 것을 어떻게 알 수 있나요?\lecturetimestamp{00:52:41}}
\lecturerresponse{측정합니다. 우주를 더 멀리 본다는 것은 더 과거를 보는 것이므로, 서로 다른 거리에 대한 관측을 통해 오늘의 값만 보는 데 그치지 않고 팽창 곡선의 과거 부분을 재구성할 수 있습니다.}
\end{classroomqa}

강의에서는 당시 이용할 수 있었던 후기 우주 관측이 대략 5퍼센트 수준에서 지수적 해와 일치했다고 보고하며, 전자 성질처럼 정밀 측정된 수준은 아니지만 증거는 설득력 있다고 설명한다. 이는 현재의 정밀도에 관한 진술이 아니라 강의 당시의 역사적 평가로 읽어야 한다.

\begin{classroomqa}
\audiencequestion{위쪽으로 휘는 팽창 곡선 때문에 결국 빅뱅이나 우주 마이크로파 배경까지 거슬러 올라가는 관측이 불가능해지나요?\lecturetimestamp{00:55:05}}
\lecturerresponse{그래프를 외삽하는 문제로는 다루지 않습니다. 여기서 중요한 불변적 문제는 사건 지평선의 존재 여부이며, 이제 살펴볼 인과적 질문이 바로 이것입니다.}
\end{classroomqa}

\section{결합계와 먼 미래}

지수 팽창이 계속되면 희석되는 모든 성분은 \(\rho_{\rm vac}\)에 비해 무시할 만큼 작아진다. 그렇다고 원자, 태양계, 은하 자체가 서로 늘어나 멀어진다는 뜻은 아니다. 서로 묶여 있을 만큼 강하게 결합된 계는 허블 흐름을 따르지 않는다.

\begin{classroomqa}
\audiencequestion{오늘날 팽창하지 않는 은하도 진공 에너지가 지배적이 되면 결국 팽창하나요?\lecturetimestamp{00:57:44}}
\lecturerresponse{결합된 계는 계속 결합된 상태로 남는다. 우리은하와 안드로메다은하는 함께 남을 근방 구조에 속한다. 반면 팽창에 이미 실려 움직이는 충분히 먼 은하들은 계속 멀어지다가 결국 우리 지평선 너머로 넘어간다.}
\end{classroomqa}

\emph{중력으로 끌리는} 것과 \emph{결합된} 것은 구별해야 한다. 탈출 속도보다 빠르게 발사된 로켓은 지구의 인력 아래에 있지만 되돌아오지는 않는다. 마찬가지로 두 은하 사이에서 보통 중력이 국소적으로 가장 큰 힘이라 해도, 두 은하의 상대 궤도가 결합되어 있다는 뜻은 아니다.

\begin{classroomqa}
\audiencequestion{우리은하는 은하군과 중력적으로 연결되어 있습니다. 가속 팽창이 결국 그 인력을 이기나요?\lecturetimestamp{00:59:17}}
\lecturerresponse{국소 구조 중 실제로 결합된 부분만 함께 남습니다. 은하단이라는 명칭만으로는 충분하지 않습니다. 결합이란 상대 운동이 탈출 속도보다 낮다는 뜻입니다. 더 먼 은하들은 서로 끌어당기면서도 탈출 속도보다 빠르게 후퇴할 수 있습니다.}
\end{classroomqa}

직관을 위해 강의에서는 상대적 방사 가속도를 인력에 해당하는 역제곱 항과 거리에 선형인 진공 항의 합으로 개략적으로 나타낸다.
\begin{equation}
\ddot r\sim-\frac{G_NM}{r^2}+H^2r.
\label{eq:l08-newtonian-vacuum}
\end{equation}
이상화된 드시터 극한에서는 계수들이 시간에 따라 변하지 않는다. 두 번째 항은 \(r\)가 커지기 때문에만 커지며, \(G_N\)가 점차 약해져서 생기는 것이 아니다.

\editorialnote{식~\eqref{eq:l08-newtonian-vacuum}은 뉴턴 역학에 따른 상대 가속도의 직관적 모형이지, 일반 상대론의 근본적인 추가 힘 법칙은 아니다. 정확한 드시터 공간에서 이웃한 자유 낙하 세계선의 상대 가속도는 \(H^2r\)에 비례한다.}

\begin{figure}[htbp]
\centering
\includegraphics[width=0.88\textwidth]{lecture_08_figure_09.png}
\caption{역제곱 인력과 거리에 선형으로 증가하는 유효 진공 반발력의 판서상 비교.}
\lectureframe{01:04:10.000}
\end{figure}

\begin{classroomqa}
\audiencequestion{바깥쪽 효과 자체가 시간에 따라 커지나요?\lecturetimestamp{01:01:18}}
\lecturerresponse{거리가 고정되어 있으면 그렇지 않습니다. 계수는 일정합니다. 자유롭게 후퇴하는 물체들은 서로 더 멀어지고, 그 거리가 커지므로 \(r\)에 비례하는 항이 더 커집니다.}
\end{classroomqa}

\begin{classroomqa}
\audiencequestion{순 효과를 뉴턴 상수가 점차 작아지는 것으로 설명할 수 있나요?\lecturetimestamp{01:02:30}}
\lecturerresponse{아닙니다. \(G_N\)는 일정합니다. 지정한 거리에서 두 계수는 모두 변하지 않으며, 거리에 따라 \(1/r^2\) 인력과 \(r\)에 비례하는 항의 상대적 중요성만 달라집니다.}
\end{classroomqa}

\begin{classroomqa}
\audiencequestion{진공 에너지가 중요해지기 전인 초기 우주에서도 같은 은하들이 중력적으로 결합되어 있었나요?\lecturetimestamp{01:05:03}}
\lecturerresponse{보통 인력이 우세했다는 사실만으로 결합을 뜻하지는 않습니다. 많은 은하가 이미 탈출 속도 정도 또는 그보다 빠른 속도로 서로 멀어지고 있었습니다. 이후 진공 에너지는 큰 규모의 팽창이 감속해 다시 수축할지, 가속해 계속 멀어질지를 결정하지만, 가장 큰 힘이 무엇인지와 궤도가 결합되어 있는지는 서로 다른 진술이다.}
\end{classroomqa}

따라서 먼 미래의 관측자는 국소 구조가 합쳐진 곳에 살면서 결합되지 않은 모든 은하를 관측할 수 없게 될 수 있다. 역사적 기록이 없다면 그 관측자는 우주가 거의 텅 비었다고 합리적으로 추론할 수 있다. 국소 천체의 파괴가 아니라 이런 인과적 고립이 지평선으로 이어지는 물리적 배경이다.

\section{사건 지평선의 의미}

휴식 시간 뒤에 다음과 같이 두고,
\begin{equation}
a(t)=e^{Ht}
\end{equation}
계량을 다음처럼 쓰자.
\begin{equation}
ds^2=-dt^2+e^{2Ht}d\mathbf{x}^2.
\label{eq:l08-de-sitter-metric}
\end{equation}

\begin{classroomqa}
\audiencequestion{이 계량이 나타내는 시공간은 무엇이라고 부르나요?\lecturetimestamp{01:08:13}}
\lecturerresponse{드시터 시공간이며, 흔히 드시터 공간이라고 줄여 부릅니다. 진공 지배 우주가 후기에는 이 기하에 가까워지므로 그 지평선이 중요합니다.}
\end{classroomqa}

인과 구조를 그리기 전에 펜로즈 도표의 역할을 떠올려 보자. 영원히 뻗는 구간을 유한한 그림 안에 압축하면서도 빛꼴 광선은 \(45^\circ\)를 유지한다. 평탄한 시공간에서 질량을 가진 물체는 미래 시간꼴 무한대에, 빛꼴 광선은 빛꼴 무한대에 도달한다. 임의로 오래 기다릴 수 있는 관성 관측자에게는 시공간의 어떤 영구적인 영역도 배제하는 사건 지평선이 없다.

\begin{figure}[htbp]
\centering
\includegraphics[width=0.80\textwidth]{lecture_08_figure_10.png}
\caption{지평선을 도입하기 전에 시간꼴 종착점과 빛꼴 종착점을 구별하기 위한 평탄 시공간의 펜로즈 도표.}
\lectureframe{01:10:35.000}
\end{figure}

블랙홀 도표를 비교 대상으로 삼아 보자. 외부 관측자는 자신의 미래 세계선의 인과적 과거에 있는 사건에서만 신호를 받을 수 있다. 관측 시간이 무한대로 가더라도 내부 영역은 여전히 배제된다. 그 관측자와 통신할 수 있는 사건과 끝내 통신할 수 없는 사건을 가르는 경계가 사건 지평선이다.

\begin{definition}
관측자의 미래 사건 지평선은 미래 방향의 인과적 신호가 언젠가 그 관측자의 세계선에 도달할 수 있는 사건들의 집합을 경계 짓는다.
\label{def:l08-event-horizon}
\end{definition}

\begin{figure}[htbp]
\centering
\includegraphics[width=0.80\textwidth]{lecture_08_figure_11.png}
\caption{블랙홀의 비교 도표: 아무리 늦은 외부 관측 시점에서 거꾸로 그은 빛원뿔도 사건 지평선 너머의 영역을 포함하지 않는다.}
\lectureframe{01:13:50.000}
\end{figure}

따라서 사건 지평선은 전역적인 인과 경계이지, 어떤 물체가 오늘 희미하게 보이는 거리만을 뜻하지 않는다. 지평선을 찾으려면 아무리 오래 기다린 뒤에도 도달할 수 있는 신호가 무엇인지 알아야 한다.

\section{무한한 미래를 압축하기}

식~\eqref{eq:l08-de-sitter-metric}의 공간 좌표는 이미 무한한 평탄 단면 전체에 걸쳐 있다. 공간 좌표는 그대로 두고 시간만 압축하자. 다음 조건을 만족하도록 등각 시간 \(T\)를 도입한다.
\begin{equation}
dt^2=e^{2Ht}dT^2.
\label{eq:l08-conformal-definition}
\end{equation}
미래 방향의 제곱근을 취하고 정리하면 다음을 얻는다.
\begin{equation}
dT=e^{-Ht}dt.
\end{equation}

\begin{classroomqa}
\audiencequestion{적분값은 부호까지 포함해 얼마인가요?\lecturetimestamp{01:17:25}}
\lecturerresponse{값은 \(-H^{-1}e^{-Ht}\)입니다. 적분 상수를 더해도 \(T\)의 원점만 이동하므로, 무한한 미래가 \(T=0\)에 놓이도록 정하면 됩니다.}
\end{classroomqa}

따라서
\begin{equation}
T=-\frac{1}{H}e^{-Ht},
\qquad -\infty<T<0.
\label{eq:l08-conformal-time}
\end{equation}
\(t\to-\infty\)일 때 \(T\to-\infty\)이다. 고유 시간 \(t\to+\infty\)일 때는 \(T\to0^{-}\)이다. 무한한 미래를 유한한 수평선 \(T=0\)에 놓았다. 어떤 관측자도 유한한 고유 시간 안에 그곳에 도달하지 않는다.

\begin{figure}[htbp]
\centering
\includegraphics[width=0.86\textwidth]{lecture_08_figure_12.png}
\caption{등각 시간 적분과 무한한 미래 \(t\to+\infty\)를 \(T=0^{-}\)에 대응시키는 관계.}
\lectureframe{01:19:10.000}
\end{figure}

식~\eqref{eq:l08-conformal-time}에서,
\begin{equation}
HT=-e^{-Ht},
\qquad
e^{2Ht}=\frac{1}{H^2T^2}.
\label{eq:l08-conformal-factor}
\end{equation}
식~\eqref{eq:l08-de-sitter-metric}에 \(dt=e^{Ht}dT\)를 사용하면 다음을 얻는다.
\begin{equation}
ds^2=\frac{1}{H^2T^2}
\left(-dT^2+d\mathbf{x}^2\right).
\label{eq:l08-conformal-metric}
\end{equation}
괄호 안의 식은 민코프스키 계량이다. 이제 드시터 기하 전체는 \(T\to0^{-}\)에서 발산하는 공통 등각 인자에 담긴다.

\begin{figure}[htbp]
\centering
\includegraphics[width=0.87\textwidth]{lecture_08_figure_13.png}
\caption{\(e^{2Ht}\)와 \(1/(H^2T^2)\)의 관계를 계산해 계량을 등각 형태로 나타내는 과정.}
\lectureframe{01:22:50.000}
\end{figure}

\begin{classroomqa}
\audiencequestion{이 계량에서 빛의 경로를 결정하는 조건은 무엇인가요?\lecturetimestamp{01:23:49}}
\lecturerresponse{빛은 영 경로 \(ds^2=0\)를 따른다. 따라서 공통 등각 인자는 소거되어, 선택한 공간 방향을 따라 \(dT=\pm dx\)가 남는다.}
\end{classroomqa}

빛은 \((T,x)\) 도표에서 직선 \(45^\circ\)으로 뻗는다.
\begin{equation}
\frac{dx}{dT}=\pm1.
\label{eq:l08-null-lines}
\end{equation}
\(T=0\)에 가까워지면 좌표 그림이 크게 압축되지만, 통신할 수 있는 사건이 무엇인지를 보여 주는 데 필요한 정보는 그대로 보존한다.

\begin{figure}[htbp]
\centering
\includegraphics[width=0.86\textwidth]{lecture_08_figure_14.png}
\caption{영 조건 \(ds^2=0\)을 적용해 \(45^\circ\) 빛살을 구하는 순간의 등각 계량.}
\lectureframe{01:24:05.000}
\end{figure}

\section{드 시터 공간의 앨리스와 밥}

앨리스를 고정된 공변 위치 \(x=0\)에 둔다. 앨리스의 세계선은 수직으로 뻗어 도표의 \(T=0\)에서 끝나지만, 앨리스의 손목시계로 잰 고유시간으로는 그곳에 도달하는 데 무한한 시간이 걸린다. 그 끝점에서 과거 방향으로 향하는 영측지선을 그리면, 앨리스가 신호를 받을 수 있는 최대 영역의 경계가 된다. 그 바깥의 사건들은 영원히 접근할 수 없다.

\begin{figure}[htbp]
\centering
\includegraphics[width=0.88\textwidth]{lecture_08_figure_15.png}
\caption{앨리스의 과거 광원뿔의 극한. 그 영측지선 경계는 관측자에 따라 달라지는 드 시터 사건지평선이다.}
\lectureframe{01:28:50.000}
\end{figure}

이 지평선은 관측자마다 다르다. 다른 관측자는 미래 끝점이 다르므로 극한 과거 광원뿔도 다르다. 한때 밥이 앨리스의 인과 영역 안에 있어 두 사람이 메시지를 주고받았다고 하자. 밥이 나중에 허블 흐름을 따라 앨리스의 지평선을 넘어가면, 그 경계를 넘은 뒤 방출된 신호는 앨리스가 아무리 오래 기다려도 도달할 수 없다. 밥에게도 자신의 지평선이 있으며, 결국 앨리스가 나중에 보내는 신호를 받을 수 있는 능력을 잃는다.

\begin{figure}[htbp]
\centering
\includegraphics[width=0.88\textwidth]{lecture_08_figure_16.png}
\caption{서로 다른 인과 경계를 가진 앨리스와 밥. 공변 관측자인 밥은 앞서 서로 신호를 주고받았더라도 앨리스의 과거 인과 영역에서 벗어날 수 있다.}
\lectureframe{01:30:20.000}
\end{figure}

앨리스와 밥이 손을 잡고 있거나, 더 현실적으로는 고유 거리를 일정하게 유지하며 서로 묶여 있다면, 두 사람 사이의 동행좌표상 간격은 \(a^{-1}\)에 따라 줄어든다. 등각도표에서는 두 사람의 세계선이 서로 가까워지지만, 그렇다고 고유 거리가 사라지는 것은 아니다. 커지는 등각인자가 줄어드는 좌표 간격을 보상하므로, 두 사람은 끝없이 통신할 수 있다.

\begin{classroomqa}
\audiencequestion{지평선이 두 사람을 갈라놓은 뒤에도 밥은 앨리스를 같은 시각에 볼 수 있나요?\lecturetimestamp{01:31:43}}
\lecturerresponse{어떤 관측자도 먼 사건을 문자 그대로 같은 시각에 보는 것은 아니다. 관측은 언제나 과거 광원뿔 안의 사건을 보는 일이다. 더 정확히 말하면, 지평선을 넘은 뒤에는 다른 세계선에서 새로 방출된 신호가 도달할 수 없다.}
\end{classroomqa}

먼 미래의 우주론적 상황에서 앨리스는 우리와 중력적으로 묶인 국소 잔해를, 밥은 각각의 중력적으로 묶이지 않은 먼 은하를 나타낸다. 은하들이 앨리스의 시계로 유한한 순간에 갑자기 사라지는 것은 아니다. 은하에서 오는 나중의 신호는 더 긴 시간 간격으로 퍼지고 더 긴 파장 쪽으로 이동하다가, 물리적으로 검출할 수 없게 된다.

\begin{classroomqa}
\audiencequestion{먼 미래에는 정말로 먼 은하들이 하나도 보이지 않는다고 이론이 말하나요?\lecturetimestamp{01:33:38}}
\lecturerresponse{중력적으로 묶이지 않은 은하들은 국소 관측자의 사건지평선 너머로 지나간다. 마지막으로 보이는 신호는 점점 더 늦게 도달하고 적색편이되어, 실용적인 의미에서는 중력적으로 묶인 국소 구조만 관측할 수 있게 된다.}
\end{classroomqa}

\section{지평선에서의 적색편이}

아주 늦은 등각시간에 밥에게서 출발한 오래된 빛 신호가 앨리스에게 도달하는 모습을 여전히 그릴 수 있으며, 그러면 앨리스가 계속 밥을 보는 듯 보일 수 있다. 물리적인 내용은 그러한 신호들의 간격과 에너지에 있다. 밥의 시계로 같은 시간 간격을 두고 출발한 신호들은 앨리스의 고유시간으로는 점점 더 큰 간격을 두고 도착한다. 진동수와 광자 에너지는 낮아진다:
\begin{equation}
\nu_{\rm received}\longrightarrow0,
\qquad
E_{\rm received}=h\nu_{\rm received}\longrightarrow0.
\end{equation}
밥의 모습은 느려지고, 붉어지며, 희미해진다. 극한에서 신호가 전달하는 유용한 정보량은 한없이 작아진다.

\begin{figure}[htbp]
\centering
\includegraphics[width=0.89\textwidth]{lecture_08_figure_17.png}
\caption{드 시터 지평선의 적색편이와 블랙홀 지평선의 익숙한 외부 관측 모습을 비교한, 판서의 마지막 상태.}
\lectureframe{01:36:20.000}
\end{figure}

\begin{classroomqa}
\audiencequestion{바깥에 있는 관측자가 블랙홀에 떨어지는 물체가 지평선을 넘는 모습을 끝내 보지 못하는 것과 비슷한가요?\lecturetimestamp{01:36:09}}
\lecturerresponse{그렇다. 두 경우 모두 관측자에게 신호는 점점 더 적색편이되어 도달하므로, 관측자는 지평선을 넘는 일이 완료되는 모습을 직접 보지 못한다. 인과 도표는 나중에 방출된 신호 중 어떤 것이 결코 도달할 수 없는지를 보여 주고, 적색편이는 마지막으로 관측 가능한 정보가 어떻게 희미해지는지를 설명한다.}
\end{classroomqa}

이 비유는 온도에 관한 질문으로 이어진다.

\begin{classroomqa}
\audiencequestion{앨리스는 지평선 바로 바깥의 온도를 어떻게 측정할 수 있을까요?\lecturetimestamp{01:39:34}}
\lecturerresponse{실제로는 온도계를 케이블에 매달아 내려 지평선을 넘지 않도록 한 뒤 다시 끌어올릴 수 있다. 드 시터 관측자도 우주론적 지평선 근처에서 이와 비슷한 실험을 생각해 볼 수 있다. 양자 이론에서 지평선은 뜨겁지만, 이번 강의는 고전적 기하학에서 멈추고 그 계산은 뒤로 미룬다.}
\end{classroomqa}

여기서는 드 시터 온도 공식을 유도하지 않는다. 호킹 복사는 이 고전적 그림에 대한 양자 보정의 예고로만 등장한다.

\section{고정된 지평선 반지름}

등각 도표를 원래 시간 좌표로 다시 표현할 수 있다. 고정된 좌표 간격에 놓인 두 공변 점의 경우
\begin{equation}
D(t)=e^{Ht}\lvert\Delta x\rvert,
\end{equation}
따라서 두 점 사이의 고유거리는 지수적으로 증가한다.

\begin{classroomqa}
\audiencequestion{원래의 변환되지 않은 시간 좌표에서는 같은 운동이 어떻게 보이나요?\lecturetimestamp{01:41:14}}
\lecturerresponse{공변 좌표가 고정되어 있으면, 공간 계량 전체가 \(e^{2Ht}\)에 따라 커진다. 따라서 고유 간격은 \(e^{Ht}\)처럼 증가한다. 좌표 변환을 되돌리면 팽창이 명시적으로 드러난다.}
\end{classroomqa}

지평선은 앨리스로부터의 공변 좌표 거리가 고정되어 있지 않으므로 다르게 움직인다. \(45^\circ\) 영측지선이 \(T=0\)에서 끝나는 것으로부터, 일정한 \(T<0\) 단면에서 지평선까지의 거리는
\begin{equation}
\Delta x_H=-T=\frac{1}{H}e^{-Ht}.
\label{eq:l08-coordinate-horizon}
\end{equation}
이다. 여기에 척도인자를 곱하면
\begin{equation}
D_H(t)=a(t)\Delta x_H
=e^{Ht}\frac{e^{-Ht}}{H}
=\frac{1}{H}.
\label{eq:l08-fixed-horizon}
\end{equation}
을 얻는다. 단위를 복원하면 \(D_H=c/H\)이다. 좌표 간격은 줄고 계량의 척도는 커지지만, 두 값의 곱은 일정하게 유지된다.

\begin{classroomqa}
\audiencequestion{시간이 지나면서 앨리스와 지평선 사이의 거리는 커지나요, 줄어드나요, 아니면 그대로인가요?\lecturetimestamp{01:42:42}}
\lecturerresponse{그대로다. 좌표 거리는 줄어들고 등각 인자는 커져서, 고유 지평선 반지름은 단위계 \(c=1\)에서 \(H^{-1}\)로 고정된다.}
\end{classroomqa}

\begin{figure}[H]
\centering
\includegraphics[width=0.90\textwidth]{lecture_08_figure_18.png}
\caption{판서의 마지막 상태: 연속된 수평 단면에서 지평선까지의 좌표 거리는 줄어들지만 고유거리는 같다.}
\lectureframe{01:43:20.000}
\end{figure}

\begin{classroomqa}
\audiencequestion{오래된 상들이 완전히 사라지지는 않으니, 앨리스는 지평선에 투영된 과거를 계속 보게 되나요?\lecturetimestamp{01:43:45}}
\lecturerresponse{고전적으로 상이 지속된다는 말은, 적색편이가 끝없이 커져 검출하기가 갈수록 어려워진다는 극한의 의미에서만 성립한다. 이 설명은 지평선과 관련된 호킹 복사를 다루지 않으며, 그 복사는 양자적 기술을 바꾸게 된다.}
\end{classroomqa}

따라서 강의는 정확한 경계에서 끝난다. 진공 에너지는 지수 팽창을 일으켰고, 지수 팽창은 관찰자에 따라 달라지는 고정된 사건의 지평선을 만들었으며, 그 지평선 부근의 신호는 주파수가 한없이 낮아졌다. 지평선의 온도와 호킹 복사는 다음에 다룰 내용이다. Susskind의 말처럼 우리가 서로의 지평선을 넘지 않는다고 가정하면, 다음 강의에서 다시 만나게 된다.